Los problemas de ruteo de vehículos ocupan un lugar central en la investigación de
operaciones, ya que articulan decisiones de asignación, secuenciamiento y transporte
bajo restricciones operacionales. Entre sus variantes, el \textit{Capacitated Vehicle
Routing Problem} (CVRP) destaca como un modelo base para estudiar la distribución
física de mercancías desde un depósito central hacia clientes con demanda conocida,
cuando la capacidad de carga de cada vehículo es limitada
\parencite{toth2002vehicle}.

\subsection{Contexto y motivación}

El interés práctico por el CVRP proviene de la necesidad de diseñar redes de
distribución con bajo costo total y alto nivel de servicio. En operaciones logísticas
reales, una planificación de rutas deficiente incrementa kilómetros recorridos,
horas-hombre, consumo de combustible y uso de flota, deteriorando la competitividad del
sistema. Desde una perspectiva de modelación, el CVRP representa la interacción entre
decisiones de partición de clientes en rutas y decisiones de ordenamiento de visitas
dentro de cada recorrido, lo que lo convierte en un problema estructuralmente complejo.

Dos contextos de aplicación ilustran esta relevancia. En distribución de retail y
comercio mayorista, un centro de distribución debe abastecer múltiples locales o puntos
de venta con una flota propia o subcontratada; aquí, el cumplimiento de capacidad y la
consolidación eficiente de carga son determinantes para reducir costos de transporte y
frecuencia de despachos. En recolección de residuos domiciliarios, los camiones deben
visitar sectores urbanos completos, respetando la capacidad de compactación antes de
retornar al depósito o estación de descarga; en este caso, la factibilidad operacional
de cada ruta depende directamente de la demanda acumulada por sector. En ambos
escenarios, el CVRP entrega una representación matemática suficientemente rica para
apoyar decisiones tácticas y operacionales.

\subsection{Fundamentos teóricos}

El CVRP pertenece a la familia de problemas $\mathcal{NP}$-hard y generaliza, entre
otros, al problema del viajante de comercio cuando la capacidad es suficientemente
grande o cuando se utiliza un solo vehículo \parencite{lenstra1981complexity}. Esta
dificultad combinatoria explica el interés por los métodos exactos basados en
programación entera, relajaciones lineales y procedimientos de ramificación, los cuales
permiten certificar optimalidad para instancias de tamaño pequeño y mediano, además de
proporcionar cotas inferiores útiles para evaluación algorítmica.

Para la lectura de este informe conviene fijar algunas nociones que articulan
dichos métodos. Un \textit{grafo no dirigido} es aquel cuyas aristas carecen de
orientación: la arista $\{i,j\}$ puede recorrerse en ambos sentidos y, en
consecuencia, $c_{ij}=c_{ji}$. Cuando la matriz de costos satisface esta
condición, el problema de ruteo se denomina \textit{simétrico}, lo que permite
formularlo sobre aristas en lugar de arcos, reduciendo a la mitad el número de
variables. Las \textit{restricciones de grado} regulan el número de aristas
activas incidentes a cada nodo: un cliente dentro de una ruta posee grado dos
(una arista por la que el vehículo llega y otra por la que sale), mientras que
el depósito, del cual parten y al cual retornan $m$ vehículos, posee grado $2m$.

Una \textit{relajación} de un problema de optimización se obtiene al eliminar o
debilitar parte de sus restricciones; como el conjunto factible resultante contiene al del problema original
(toda solución factible del problema original sigue siéndolo en la
relajación), el valor óptimo de la relajación constituye una \textit{cota
inferior} del valor óptimo verdadero, es decir, un piso
garantizado bajo el cual éste no puede encontrarse. En particular, la
\textit{relajación lineal} de un programa entero se obtiene al sustituir la
condición de integralidad por su versión continua (por ejemplo,
$x_{ij}\in\{0,1\}$ por $0 \le x_{ij} \le 1$). De manera complementaria, toda solución factible del problema proporciona una
\textit{cota superior}, y la brecha relativa entre ambas cotas (\textit{gap})
acota la distancia máxima entre la mejor solución disponible y el óptimo: un
gap de $g\%$ certifica que dicha solución dista a lo más un $g\%$ del valor
óptimo, aun cuando éste no se conozca. Las cotas cumplen así un rol doble:
cuando la brecha se anula, certifican optimalidad, y durante la búsqueda
habilitan la poda en \textit{Branch and Bound}, pues toda rama cuya cota
inferior excede el costo de la mejor solución entera conocida puede descartarse
sin explorarla.

En este informe se adopta el enfoque exacto clásico para el CVRP simétrico
sobre un grafo no dirigido, basado en una estrategia de \textit{Branch and
Bound} en la que las restricciones exponenciales de subtour y capacidad no se
imponen desde el inicio \parencite{laporte1983branch}. En su lugar, el problema
se resuelve primero sobre una relajación con restricciones de grado, y las
desigualdades pertinentes se agregan únicamente cuando una solución entera
candidata presenta subtours ilegales o incumplimientos de capacidad. Este
esquema combina tres ideas fundamentales de la optimización exacta: relajación
para obtener cotas inferiores, separación de restricciones para fortalecer la
formulación, y ramificación para recuperar integralidad.

\subsection{Objetivos}

El objetivo general de este informe es resolver el CVRP mediante métodos exactos
clásicos, implementando computacionalmente la formulación de
\textcite{laporte1983branch} para el caso simétrico con generación iterativa de
desigualdades de capacidad, y evaluando el procedimiento sobre instancias
estándar de la literatura.

Los objetivos específicos son los siguientes:
\begin{enumerate}
    \item Definir el CVRP en lenguaje natural e identificar contextos logísticos
          representativos.
    \item Presentar una formulación matemática compacta del problema e interpretar sus
          restricciones principales.
    \item Revisar el estado del arte de los métodos exactos aplicables al CVRP,
          posicionando el enfoque seleccionado dentro del panorama de algoritmos
          \textit{Branch and X}.
    \item Implementar el procedimiento de resolución exacta, basado en la
          formulación de dos índices con generación iterativa de desigualdades de
          capacidad redondeada (RCI), delegando la resolución de cada modelo entero
          al solver COIN-OR CBC.
    \item Evaluar el método sobre instancias de los Sets E y A de CVRPLIB, contrastando las
          soluciones obtenidas con los valores óptimos de referencia y analizando el
          esfuerzo computacional en función del tamaño de la instancia.
\end{enumerate}

\subsection{Organización del informe}

El Capítulo~2 resume el estado del arte de los métodos exactos para el CVRP,
posicionando el enfoque adoptado dentro de la familia \textit{Branch and X}. El
Capítulo~3 presenta la definición del problema, la formulación matemática completa y
el procedimiento de resolución basado en \textit{Branch and Bound} con generación
iterativa de desigualdades de capacidad. El Capítulo~4 reporta los resultados de la experimentación computacional sobre
instancias de los Sets E y A de CVRPLIB y su análisis. El Capítulo~5 discute
los hallazgos en relación con los objetivos del trabajo, sus limitaciones y las
mejoras posibles. Finalmente, el Capítulo~6 presenta las conclusiones.